\documentclass[a4paper,11pt]{article}
\usepackage[margin=27mm]{geometry}
\usepackage{fontspec,amsmath,amssymb,amsthm,unicode-math,booktabs}
\setmainfont{STIXGeneral}\setmathfont{STIX Two Math}
\newtheorem{theorem}{Theorem}\newtheorem{lemma}[theorem]{Lemma}\newtheorem{proposition}[theorem]{Proposition}
\newcommand{\gen}[1]{\langle#1\rangle}\newcommand{\Aut}{\operatorname{Aut}}
\title{Regular thin covers with disconnected face intersections}
\author{Human Compression project}\date{Author edition; pending independent review}
\begin{document}\maketitle
\begin{abstract}
We express the cyclic voltage family of Leemans and Toledo as a fiber product
with a dihedral group. Two independent binary modules describe its endpoint
parabolics. Their intersections then follow from the same projection and kernel
argument used for mixes. Taking a canonical double cover gives the faithful,
thin, nonpolytopal family; adjoining commuting colors gives every higher rank.
An isolated finite presentation certificate identifies the exact source base;
the uniform conclusions are proved in writing. We also give an explicit
counterexample to the source's thin-base auxiliary lemma and recover its
faithful-base canonical-double-cover result with the repaired common theorem.
\end{abstract}

\section{Two modules over a polygon group}
Let $E=\mathbb F_2^2$, with basis $e=(1,0)$ and $f=(0,1)$, and put
\[
B=\begin{pmatrix}1&0\\1&1\end{pmatrix},\qquad
C=\begin{pmatrix}1&1\\0&1\end{pmatrix}.
\]
They are involutions, $Ce=e$, $Bf=f$, and $BC$ has order~3. On the three
nonzero vectors they generate the full symmetric group. For a positive
integer $m$, let
\[
D_m=\gen{b,c\mid b^2=c^2=(bc)^{3m}=1}
\]
be the dihedral group of order $6m$. Its action on $E$ sends $b$ to $B$
and $c$ to $C$, and factors through the quotient $D_m\to D_1\cong S_3$.
Use this action on both summands of $V=E_0\oplus E_3$, and form
$G_m=V\rtimes D_m$. Multiplication is
$(v,d)(w,h)=(v+d(w),dh)$. Set
\[
s_0=(e,0;1),\quad s_1=(0,0;b),\quad
s_2=(0,0;c),\quad s_3=(0,f;1).
\]
Each mark is an involution. They generate $G_m$: $s_1,s_2$ generate
$D_m$, while their conjugates of $s_0$ and $s_3$ span the two independent
copies of $E$. Nonconsecutive marks commute, because $Ce=e$, $Bf=f$,
and translations commute. Moreover
\[
|s_0s_1|=4,\quad |s_1s_2|=3m,\quad |s_2s_3|=4.       \tag{1}
\]
For example $(s_0s_1)^2=(e+Be,0;1)=(f,0;1)\ne1$;
the other outer calculation is dual.

\begin{lemma}\label{ends}
The indicated marking of $G_m$ is a string C-group. Its order is $96m$,
its endpoint parabolics are $E_0\rtimes D_m$ and $E_3\rtimes D_m$, and
its face numbers are $(4,6m,6m,4)$.
\end{lemma}
\begin{proof}
The two endpoint descriptions follow from the spanning argument above.
Their intersection is precisely the common complement $D_m$, because
$E_0\cap E_3=0$. This is a direct application of the common criterion using image and
kernel: their images in $D_m$ coincide, and their translation
kernels are disjoint.

Inside the first endpoint, $\gen{s_0,s_1}=E_0\rtimes\gen{b}$ has order~8,
and $\gen{s_1,s_2}=D_m$. Their intersection is exactly $\gen{b}$.
Both are dihedral rank-two C-groups, so the rank-three endpoint has the
intersection property. To make the inference precise, apply the recursive
path criterion from \emph{Intersections in simultaneous group actions},
Lemma~2, to its regular action: both polygon ends are polytopal and their
intersection is the common marked subgroup. The dual argument proves
the other endpoint, using
$\gen{s_2,s_3}=E_3\rtimes\gen{c}$ and intersection $\gen{c}$ with $D_m$.
The same recursive criterion in rank~4, together with the endpoint
intersection $D_m$, proves the full C property, including every subset
of the four marks. It is not inferred just from adjacent orders.

The endpoint orders are $24m$, and $|G_m|=16(6m)=96m$.
Deleting color~1 leaves
$\gen{s_0}\times\gen{s_2,s_3}\cong C_2\times D_8$, of order~16;
deleting color~2 is dual. Indices give the four displayed face counts.
\end{proof}

\begin{lemma}\label{odd}
Both endpoint flag graphs are nonbipartite, for every $m$.
\end{lemma}
\begin{proof}
The three nonzero vectors in $E_0$ sum to zero. Conjugating $s_0$ by
suitable words in $b,c$ realizes these three translations. Their product
is the identity. Written in the four marked involutions, this relation
has odd length: each conjugated translation contributes one occurrence
of $s_0$ and twice the length of its conjugating word. It is an odd
closed walk in the first endpoint. The same argument in $E_3$ proves
the second. This proves nonbipartiteness for all $m$; it does not depend
on an order calculation in sample covers.
\end{proof}

\section{The exact voltage construction}
Write $P_m$ for the regular polytope just proved. First consider $m=1$.
Its endpoint groups $E\rtimes S_3$ act faithfully as all 24 affine
permutations of the four-element set $E$, and hence are $S_4$. The
products $s_0s_1s_2$ and $s_1s_2s_3$ have order~3: the projected rotation
has order~3 and its norm on $E$ is $1+BC+(BC)^2=0$. Thus the facet and
vertex-figure types are respectively $\{4,3\}_3$ and $\{3,4\}_3$.
The exact universal presentation with these relations has order~96,
as already certified in the accepted common projective marking paper's
sixth universal row. The complete isolated program \texttt{base\_certificate.g} verifies that
exact presentation and its marked isomorphism to this 16-point affine
permutation model; it is a required finite premise for this source
identification, and its original attribution is retained. Since the displayed
$G_1$ has order~96 and is a marked quotient of that presentation, this
identifies it with the source base
$\mathcal B=\{\{4,3\}_3,\{3,4\}_3\}$.

Use right multiplication for the colored Cayley adjacencies of $G_1$.
Let $V=E_0\oplus E_3$ denote its translation subgroup. Assign integer
voltages to color~2 darts by
\[
z(g,gs_2)=\begin{cases}1&g\in V,\\-1&g\in Vs_2,\\0&\text{otherwise},\end{cases}
\]
and voltage0 to all other colors. The assignments on reversed darts
are negatives, since $s_2^2=1$. Reduce modulo $m$ for the cyclic cover.

\begin{proposition}\label{sourcevoltage}
This is, up to reversing all voltages and choosing the opposite edge
pair, the voltage assignment of Leemans--Toledo Section~5. Its cover is
the Cayley maniplex of $G_m$ with the four marks above.
\end{proposition}
\begin{proof}
An edge face of $P_1$ through the identity is
$J=\gen{s_0,s_2,s_3}$, of order~16. Its intersection with $V$ has order~8,
and $VJ=V\rtimes\gen{c}$ has order~32. Hence $VJ$ is the union of two
edge faces. Vertices of $P_1$ identify with the four vectors of $E_0$;
the endpoints of $J$ are $0,e$, and those of the other edge face are
the other two vectors. These are exactly an opposite edge pair in the
source's $K_4$ skeleton.

The 16 origins in $V$ are closed under colors~0 and3, and their color~1
and color~2 neighbors lie outside $V$. They are also closed under the
word101, since $s_1s_0s_1$ is a translation. Thus they satisfy all three
orientation rules accompanying source Figure~5.1. These rules determine
the orientation once one origin is selected. Indeed the moves
$s_0,s_3,s_1s_0s_1$ preserve the sign, and $s_2$ reverses it.
These moves generate $VJ$: the first and third span $E_0$, while
$s_3$ and its conjugate by $s_2$ span $E_3$. Thus the 32 flags of the
opposite edge pair have their signs uniquely determined. The two
choices are $V$ and $Vs_2$, as asserted.

The quotient of the base flag graph by $V$ is the Cayley six-cycle of
$D_1$ in colors~1,2, with semiedges in colors~0,3. Our voltage is pulled
back from a voltage on that six-cycle: one directed color~2 edge has
voltage1 and the other five edges have voltage0. Its derived cover is
the alternating cycle of length $6m$, the Cayley graph of $D_m$ in
$b,c$. Its projection to the six-cycle is the marked quotient
$D_m\to D_1$.

Pulling this cover back to the base gives the pairs
\[
\{(g,d)\in G_1\times D_m: gV=\overline d\in D_1\}.
\]
The explicit map $(v,d)\mapsto((v,\overline d),d)$ is a group
isomorphism from $V\rtimes D_m$ onto this fiber product, with inverse
$((v,\overline d),d)\mapsto(v,d)$. It sends the four displayed marks
to the four colored adjacency pairs. The subgroup generated by them
is the entire product, since it contains the complement $D_m$ and
both translation summands. Thus its Cayley graph is precisely the
connected voltage cover, rather than merely one unidentified component.
This proves the assertion for every modulus simultaneously.
\end{proof}

The kernel of $G_m\to G_1$ is $\gen{(bc)^3}$, of order $m$.
It is contained in both endpoint parabolics. Their full preimages are
therefore connected endpoint faces. Each middle parabolic has order16
and maps isomorphically to the corresponding order16 base parabolic;
its intersection with the kernel is trivial. Every base middle face
therefore has $m$ isomorphic lifts. This proves exactly the face-preimage
clauses of source Lemma~5.1, rather than just their eventual consequences.

The source proves regularity by lifting four base automorphisms,
including a piecewise lift with nine adjacency cases. The actual fiber
product description makes those cases unnecessary: left multiplication
in $G_m$ is already regular on its Cayley graph. Lemma~\ref{ends}
also proves every intersection needed for polytopality directly. The
retained source construction and its attribution are unchanged; this
is a different complete derivation of that construction's properties.

For completeness, the source's displayed exceptional lift is also recovered.
Choose the identity base flag in $V$, so $\rho_i$ is left multiplication
by $s_i$. Left multiplication by $s_0,s_3$ preserves the voltage, and left
multiplication by $s_2$ negates it. For $\rho_1$, put $h(q)=0$ on
$q\in\{1,b\}\subset D_1$ and $h(q)=1$ on the other four quotient flags.
The lift is
\[
(g,j)\longmapsto(s_1g,-j+h(gV)).                 \tag{2}
\]
The following complete six-state table verifies the only nonzero-voltage
color. Its entries are integer equalities, so apply modulo every $m$.
\[
\begin{array}{c|rrrrrr}
q&1&b&c&bc&cb&bcb\\\hline
z(q,qc)&1&0&-1&0&0&0\\
z(bq,bqc)&0&1&0&-1&0&0\\
h(qc)-h(q)&1&1&-1&-1&0&0
\end{array}
\]
The second row is the negative of the first plus the third. Colors~0,3
do not change $q$, and right multiplication by $b$ preserves $h$, so
the required identity for all other colors is zero on both sides.
The common lift calculation proves (2). Since
$V\cup s_1V$ is exactly the source set $F_X\cup F_X^{\rho_1}$,
this is its literal piecewise formula. A deck translation adjusts any
of the other lifted base generators to the prescribed adjacent flag.
Thus the lift statement and its concrete formula are retained, even
though they are unnecessary as separate premises of the regularity proof.

\section{The nonpolytopal family and all higher ranks}
\begin{theorem}\label{main}
For each positive integer $m$, the canonical double cover of $P_m$ is
a finite regular faithful thin nonpolytopal rank-four maniplex with
$192m$ flags and face numbers $(4,12m,12m,4)$.
\end{theorem}
\begin{proof}
Lemmas~\ref{ends} and~\ref{odd} verify every hypothesis of the common
canonical-double-cover corollary in \emph{Intersections and faces in
graph covers}. Its proof gives all conclusions. In particular the
vertex--facet intersections split into two components; faithfulness and
thinness do not imply the missing strong connectivity. Distinct values
of $m$ give distinct flag counts and therefore nonisomorphic maniplexes.
\end{proof}

For clarity the failure also has a group description. Nonorientability
makes the double-cover group the entire $G_m\times C_2$, with marks
$(s_i,1)$. Both endpoint parabolics contain the deck involution
$(1,1)$, because each has an odd relation. Their intersection is
$D_m\times C_2$. The middle subgroup generated by $(b,1),(c,1)$ is
the graph of the usual word-length parity homomorphism on the dihedral
group, and has order $6m$. The deck involution is consequently an
explicit element of the endpoint intersection outside the middle
subgroup. This is the common fiber obstruction, not a failure detected
only through a finite graph search.

Let $M$ be any rank-$r$ maniplex. Define $R(M)$ on
$\operatorname{Flags}(M)\times\{0,1\}$ by retaining all old adjacencies
in each copy and adjoining color $r$ which interchanges the copies.
Its old $i$-faces are $F_i\times\{0,1\}$ and its two new facets are
$M\times\{0\}$ and $M\times\{1\}$. These descriptions show directly
that faithfulness and thinness hold exactly when they do in $M$:
maximal chains choose one new facet and then one old maximal chain;
all old diamond conditions are unchanged, and the new top diamond has
exactly the two new facets. A face-chain intersection either fixes one
copy, or is the Cartesian product of its old intersection with a single
edge. It is connected exactly when the old intersection is. Thus the connected-intersection property is preserved and reflected.
For the faithful maniplexes used here it is equivalent to polytopality.
The face description also proves the source's poset-polytopality equivalence
without a faithfulness restriction. Maximal chains are pairs consisting of
an old chain and one new facet. Within one new facet, strong connectivity
is exactly that of the old poset. For chains in different new facets,
first connect the old chains while fixing their common old faces, then
change just the new facet. Conversely, applying strong connectivity to
two chains in the same new facet forces a connection in that facet and
hence in the old poset. Combined with the diamond calculation, this proves
that the new face poset is a polytope exactly when the old one is. Every automorphism acts by the same old automorphism
on both copies, optionally interchanging them, so
$\Aut(R(M))=\Aut(M)\times C_2$.
Iterating this construction $r-4$ times in Theorem~\ref{main} gives
$192m\,2^{r-4}$ flags and proves the complete principal assertion of
Leemans--Toledo Theorem~1.1 for every $r\geq4$.

\section{A thin base for which the auxiliary lemma fails}
The general thin-base Lemma~3.5 in the source is false. All coordinates
in the following example belong to $\mathbb F_2$. Start with the32 flags
$(v_0,v_1,v_2,v_3,z)$. The four adjacencies flip their indicated $v_i$;
in addition color~0 adds $v_1$ to $z$ and color~2 adds $v_3$ to $z$.
Colors~1,3 leave $z$ fixed. These are distinct fixed-point-free involutions,
and nonconsecutive pairs commute. The graph is connected: the $v_i$
can be changed freely, and the commutator of colors~0,1 changes $z$
without changing them. The same is true of the commutator of colors~2,3.

Its $i$-faces are exactly the sets with fixed $v_i$. Indeed the remaining
three colors change the other three $v$ coordinates, and at least one
of those two commutators remains available to change $z$. Thus every
face has16 flags. Every pair of faces of distinct ranks meets, so the
face poset has two elements at each proper rank, all comparable across
different ranks. It is thin. It is unfaithful: fixing all four faces
leaves both possible values of $z$.

Take a twofold voltage cover by adjoining $\epsilon\in\mathbb F_2$.
Colors~0,3 have voltage0, color~1 has voltage $v_2$, and color~2 has
voltage $v_3$. Equivalently the four permutations are
\[
\begin{array}{c|c|c|c}
\text{color}&\text{coordinate flipped}&z\text{ increment}&\epsilon\text{ increment}\\\hline
0&v_0&v_1&0\\
1&v_1&0&v_2\\
2&v_2&v_3&v_3\\
3&v_3&0&0
\end{array}                                                   \tag{3}
\]
Each voltage is unchanged on reversing its edge. Directly from (3) the
maps are distinct fixed-point-free involutions and the distant pairs
commute. A color-$1,2$ commutator changes only $\epsilon$; the
color-$2,3$ commutator changes both $z$ and $\epsilon$. Hence the preimage
of each $0$-face is connected. Dually colors~0,1 change only $z$ and
colors~1,2 change only $\epsilon$, so each $3$-face preimage is connected.
These facts also show the entire cover connected.

On a $1$-face the voltage is the difference of the potential
$z+v_0v_1$: this holds along each of its allowed colors~0,2,3 by (3).
Its two lifts are therefore indexed by
\[
(v_1,\alpha),\qquad \alpha=\epsilon+z+v_0v_1.
\]
They each map isomorphically to that base face, by the elementary
voltage-path argument. On a $2$-face the potential is $v_1v_2$, and
its two isomorphic lifts are indexed by
\[
(v_2,\beta),\qquad \beta=\epsilon+v_1v_2.
\]
All face-preimage hypotheses of the printed lemma are now verified.
The cover has face counts $(2,4,4,2)$ and is faithful: its faces determine
all four $v_i$, then $\epsilon$ from $\beta$, and $z$ from $\alpha$.

Nevertheless it is not thin. Fix the $0$-face $v_0=0$ and the $2$-face
$(v_2,\beta)=(0,0)$. For every $(a,d)\in\mathbb F_2^2$, the flag
\[
(v_0,v_1,v_2,v_3,z,\epsilon)=(0,a,0,0,d,0)
\]
lies in both these faces and in the $1$-face $(v_1,\alpha)=(a,d)$.
There are four intervening $1$-faces, rather than two. This is a direct
written counterexample to the exact thin-base statement. The independent
finite diagnostic verifies all faces and diamond intervals, but is not
a premise of this argument. The source proof's inference that an
intersection of preimages has precisely the same component count as
one preimage is not valid without a connectivity argument. The common theorem for
faithful thin bases supplies that missing argument and proves all
canonical-double-cover conclusions on their actual faithful hypothesis.

\section{Attribution and the remaining source boundary}
We use Leemans and Toledo's \emph{Faithful and thin non-polytopal
maniplexes}.

The fixed edition is arXiv:2305.01372v2, dated 10 May 2023,
with 12 physical pages. Its main finite family, the voltage choice, the small-rank
obstruction and the raviolo construction are attributed to that source;
no historical novelty is asserted for our alternative derivation.
The common paper includes the complete rank-three argument, the general
faithful middle-face lemma and every clause of Lemma~3.8. Proposition~5.4
correctly prints ``without the CIP'', exactly as proved here.

The 11-cell and 57-cell examples follow from the same common corollary.
Their retained local types are hemi-icosahedral facets with hemi-dodecahedral vertex-figures for
the 11-cell, and the dual pair for the 57-cell. These projective maps are nonorientable,
so their flag graphs have odd closed walks. The existence and regular
polytopality of these objects of Gr\"unbaum and Coxeter are the already accepted
proper-end classifications in the projective-marking library; they are
retained inputs for these illustrations, not new classifications here.
Their double covers are consequently regular, faithful, thin and
nonpolytopal. The informal exhaustive claim on source page 7 that the
192-flag example is the smallest regular example is a separate ancillary
census and is not reproduced. Nor do we re-establish the closing atlas
counts of regular 11-cell covers, construct the proposed further voltage
families, or prove the omitted four-facet extension statement.

The source's broad background assertion that $\operatorname{Pos}$ and
$\operatorname{Man}$ are inverse for every thin flagged poset is not used.
Our chain and face arguments start with maniplexes, and our inverse
identifications with polytope flag graphs have their faithful and thin
hypotheses explicitly proved. The source's graph-lifting attribution is
retained in the common sufficient lift lemma. The cyclic cover here is
an actual fiber product with a proved inverse map and proved generation;
this is where the common theory replaces repeated local intersection
and lift arguments. The finite base-presentation certificate, the
written uniform proof, and bounded validation programs remain
separate obligations. Four new Lean declarations prove only the generic
equivariant-projection intersection criterion and the actual voltage-lift
bijection and adjacency calculation. Two existing declarations concern
simultaneous action fibers. Neither file formalizes face connectivity,
the finite presentation calculation, or the uniform family; those
application obligations are the written arguments above.
\end{document}
